\section*{Abstract}
In this manuscript, I present the work done during my PhD on confined Brownian motion. Brownian motion is the erratic movement of microscopic particles when immersed into a fluid. Thanks to Einstein and followers, it is generally possible to describe Brownian motion using simple equations. However, in the last two decades a scientific revolution has taken place with the advent of miniaturization and in particular microfluidics, thanks to which we are now able to create complex networks of pipes at the micrometer scale. Microfluidics makes it possible to sort particles, such as drops, cells or bubbles, but also to distribute drugs in cells and observe their effect on thousands of them. Regarding Brownian motion, it has been observed that once confined near a wall, a particle moves much slower due to non-slip boundary conditions at the wall. The mobility is thus modified by confinement-induced effects.

My thesis work consists in experimentally measuring, analyzing and modeling the movement of micrometric colloids diffusing near a wall. To track the motion of confined Brownian microparticles, I use Lorenz-Mie holography. The Lorenz-Mie framework allows me to record the thermally-induced tridimensional trajectories of individual microparticles, within salty aqueous solutions, in the vicinity of a rigid wall, and in the presence of surface charge with a nanometric resolution. From the recorded trajectory, I construct the time-dependent position and displacement probability density functions, and study the non-Gaussian character of the latter which is a direct signature of the hindered mobility near the wall. Based on these distributions, I implement a novel, robust and self-calibrated multifitting method, allowing for the thermal-noise-limited inference of diffusion coefficients spatially-resolved at the nanoscale, equilibrium potentials, and forces at the femtonewton resolution. Moreover, I use this novel tool in order to infer non-conservative forces and study long-time higher-order statistical properties. Our goal for the future is to use this novel tool for the in-depth study of various problems relevant to nanophysics and microbiology.
In chapters 1 and 2, I present a general introduction of the subject (including en French translation). In chapter \ref{sec:chapter1}, I present the history of Brownian motion, starting from how it has been discovered to its numerical simulation. In chapter \ref{sec:chapter2}, I present state-of-the-art methods that permit tracking single colloids. In particular, I detail how we use the Lorenz-Mie framework. Then, I show the viability of this framework by studying a well-known problem: the sedimentation of a colloid. In chapter \ref{chap3}, I present the confined Brownian motion theory and its experimental study. In particular, I explain  the implementation of a novel, robust and self-calibrated multifitting method. In chapter \ref{chapfin}, I present other applications the newly developed method. Finally, I conclude the manuscript in chapters 7 and 8 (including en French translation).

\newpage

\section*{Résumé}

Dans ce manuscrit, je présente le travail effectué pendant mon doctorat sur le mouvement brownien confiné. Le mouvement brownien est le mouvement erratique de particules microscopiques lorsqu'elles sont immergées dans un fluide. Grâce à Einstein et ses successeurs, il est généralement possible de décrire le mouvement brownien à l'aide d'équations simples. Cependant, au cours des deux dernières décennies, une révolution scientifique a eu lieu avec l'avènement de la miniaturisation et en particulier de la microfluidique, grâce à laquelle nous sommes désormais en mesure de créer des réseaux complexes de tuyaux à l'échelle micrométrique. La microfluidique permet de trier des particules, comme des gouttes, des cellules ou des bulles, mais aussi de distribuer des médicaments dans des cellules et d'observer leurs effets sur des milliers d'entre elles. En ce qui concerne le mouvement brownien, il a été observé qu'une fois confinée près d'une paroi, une particule se déplace beaucoup plus lentement en raison des conditions limites non glissantes à la paroi. La mobilité est donc modifiée par les effets induits par le confinement.

Mon travail de thèse consiste à mesurer, analyser et modéliser expérimentalement le mouvement de colloïdes micrométriques diffusant près d'une paroi. Pour suivre le mouvement de microparticules browniennes confinées, j'utilise l'holographie de Lorenz-Mie. Le cadre de Lorenz-Mie me permet d'enregistrer les trajectoires tridimensionnelles thermiquement induites de microparticules individuelles, dans des solutions aqueuses salées, à proximité d'une paroi rigide, et en présence d'une charge de surface avec une résolution nanométrique. À partir de la trajectoire enregistrée, je construis les fonctions de densité de probabilité de position et de déplacement en fonction du temps, et j'étudie le caractère non-gaussien de ces dernières qui est une signature directe de la modification de la mobilité à proximité de la paroi. Sur la base de ces distributions, j'implémente une nouvelle méthode de multifitting robuste et auto-calibrée, permettant l'inférence limitée par le bruit thermique des coefficients de diffusion résolus spatialement à l'échelle nanométrique, des potentiels d'équilibre et des forces à la résolution du femtonewton. De plus, j'utilise ce nouvel outil pour déduire les forces non-conservatives et étudier les propriétés statistiques d'ordre supérieur à long terme. Notre objectif pour l'avenir est d'utiliser ce nouvel outil pour l'étude approfondie de divers problèmes liés à la nanophysique et à la microbiologie.
Dans les chapitres 1 et 2, je présente une introduction générale du sujet (y compris la traduction française). Dans le chapitre \ref{sec:chapter1}, je présente l'histoire du mouvement brownien, de sa découverte à sa simulation numérique. Dans le chapitre \ref{sec:chapter2}, je présente les méthodes les plus récentes qui permettent de suivre des colloïdes uniques. En particulier, je détaille comment nous utilisons le cadre de Lorenz-Mie. Ensuite, je montre la viabilité de ce cadre en étudiant un problème bien connu : la sédimentation d'un colloïde. Dans le chapitre \ref{chap3}, je présente la théorie du mouvement brownien confiné et son étude expérimentale. En particulier, j'explique l'implémentation d'une nouvelle méthode de multifitting, robuste et auto-calibrée. Dans le chapitre \ref{chapfin}, je présente d'autres applications de la méthode nouvellement développée. Enfin, je conclus le manuscrit dans les chapitres 7 et 8 (y compris la traduction française).

\vspace{4cm}

\begin{figure}[h]
	\begin{center}
		\includegraphics{01_head/images/RENDU_16_10_MAX.png}
		\caption{Artist's view of the Lorenz-Mie method. A polystyrene ball (in white) is illuminated from above. The interference pattern is indicated by the concentric rings (Credit: \href{https://www.behance.net/pierresavary357}{Pierre Savary}). \\ Vue d'artiste de la méthode de Lorenz-Mie. Une bille de polystyrène (en blanc) est éclairée par le haut. La figure d'interférences est indiquée par les anneaux concentriques (Crédit : \href{https://www.behance.net/pierresavary357}{Pierre Savary}).}
	\end{center}
\end{figure}

